\section{Primera empresa: de las herramientas para WeChat a las lecciones de comercialización}

La sección anterior trazó la ruta; esta sección examina primero la primera empresa de Xiao Hong (肖弘). En 2015, al egresar de la Universidad de Ciencia y Tecnología de Huazhong, emprendió de inmediato. Al principio desarrolló productos para el entorno universitario que tuvieron mala respuesta; después pasó a un editor para cuentas oficiales de WeChat y, más tarde, a un SCRM para WeCom (企业微信). La formación más importante de esta etapa no fue solo haber creado productos, sino aprender del fracaso, del flujo de efectivo, de la comercialización, del capital y de la venta de la empresa: en un producto de herramienta, una historia más grande no es mejor; lo que hace falta es encontrar usuarios reales, pagos reales y una economía unitaria sostenible.

\begin{figure}[H]
\centering
\includegraphics[width=0.92\textwidth]{founder-cycle.png}
\caption{Ciclo del emprendedor serial: en habilidades, capital, mercado, producto y mentalidad aparecen continuamente nuevas dimensiones. Diagrama conceptual propio, elaborado a partir de la conversación entre 00:05:15 y 00:58:15.}
\end{figure}

\begin{knowledgebox}{Lectura de la figura: la capacidad emprendedora no se completa de una sola vez}
Alrededor del Founder, la figura muestra los puntos de habilidad, el capital, el mercado, el producto, la mentalidad y la negociación estratégica. Xiao Hong repite que cada vez que sentía que sus puntos de habilidad ya eran suficientes, surgía una dimensión nueva. En la primera empresa aprendió primero de producto; después, de comercialización, y más adelante descubrió que el capital también era una dimensión nueva.
\end{knowledgebox}

\subsection{De la retroalimentación positiva en la universidad a la negativa en el mundo real}

Esta subsección explica los primeros reveses. En la universidad, Xiao Hong participó en un club técnico, en un centro de nuevos medios, en un blog y en proyectos de producto, y recibió mucha retroalimentación positiva. Al egresar, desarrolló productos como una red social anónima y un mercado de artículos de segunda mano, pero no consiguieron usuarios, y el equipo estuvo a punto de irse a trabajar a Pekín. Esta experiencia muestra que los pequeños éxitos en el entorno universitario no pueden extrapolarse directamente al mercado real. Emprender de verdad es más duro, porque los usuarios no son el círculo de conocidos, los canales no existen de forma natural y el pago no ocurre por sí solo.

\begin{warningbox}{Error común: tomar la retroalimentación positiva de la universidad como validación de mercado}
Crear herramientas, cuentas oficiales o pequeños productos en la universidad demuestra capacidad e interés personales, pero no demuestra que exista un mercado comercial. La primera lección después de emprender suele ser descubrir que lo que se había recibido era retroalimentación positiva local, y no una demanda amplia.
\end{warningbox}

\subsection{Por qué la comercialización debe empezar antes}

Esta subsección examina la comercialización del editor para cuentas oficiales. En su retrospectiva, Xiao Hong dice que el equipo dedicó demasiado tiempo al crecimiento de usuarios y a la iteración del producto, y pensó poco en la comercialización. Más tarde decidieron vender software, como un SaaS, en lugar de crear una plataforma de intermediación publicitaria. La intermediación publicitaria parece tener un techo más alto, pero no necesariamente forma parte de las capacidades centrales de un equipo de software; vender software parece más pequeño, pero es más fiel a lo que el equipo sabe hacer, cierra mejor el ciclo y además permite al equipo mejorar el producto de forma continua.

\begin{importantbox}{Experiencia práctica: un techo alto no significa que se pueda lograr}
Muchos equipos emprendedores, para contar una historia más grande, eligen modelos de negocio que no corresponden a sus capacidades. En el caso del editor para cuentas oficiales, «captar clientes con la herramienta y después hacer intermediación publicitaria» sonaba más grande, pero lo que realmente funcionó fue el software por suscripción. En un producto de herramienta, que los usuarios paguen directamente por el software forma, de hecho, un ciclo cerrado más sano.
\end{importantbox}

\subsection{El VC es un medio costoso de obtener fondos}

Esta subsección entra en el tema del capital. Xiao Hong dice que lo costoso del VC no se nota cuando a la empresa le va mal, sino cuando le va bien: si la empresa tiene un gran éxito, el costo de la financiación temprana con capital accionario se multiplica. Financiarse es una herramienta razonable, pero el fundador debe saber que cambiará la historia, los objetivos y la presión. Recaudar demasiado dinero puede obligar a la empresa a contar una historia más grande y a hacer cosas que no necesariamente son correctas.

\begin{figure}[H]
\centering
\includegraphics[width=0.90\textwidth]{vc-expensive-capital.png}
\caption{El VC es un medio costoso de obtener fondos: lo costoso no aparece en los momentos difíciles, sino cuando a la empresa le va bien. Diagrama conceptual propio, elaborado a partir de la conversación entre 00:33:16 y 00:42:10.}
\end{figure}

\begin{knowledgebox}{Lectura de la figura: el capital es una herramienta, no una dirección}
A la izquierda, en los momentos difíciles, la financiación parece dinero de rescate; a la derecha, cuando a la empresa le va bien, la dilución temprana y los compromisos mayores se multiplican. Al leer esta figura conviene retener lo siguiente: el VC no es malo; lo esencial es que el Founder sepa qué restricciones obtiene a cambio de usar esta herramienta.
\end{knowledgebox}

\subsection{El inversionista no es el jefe}

Esta subsección conserva una founder voice de mucho valor. Xiao Hong recomienda no tratar a los inversionistas como jefes: si un inversionista comparte un artículo, ese gesto debería pesar más o menos como la opinión de un KOL, y no convertirse automáticamente en estrategia de la empresa. Si un inversionista de verdad tiene una opinión firme, debe convocarse una reunión formal y dejar registro de quién sostiene qué juicio, en lugar de dejar que el Founder adivine si el inversionista quiere que la empresa cambie de rumbo. La razón es que el apalancamiento del Founder es altísimo: un solo error de juicio arrastra el dinero y la dirección del equipo.

\begin{importantbox}{Nota de clase: las decisiones de alto apalancamiento requieren un mecanismo formal}
El Founder puede escuchar la opinión de los inversionistas, pero no debe tomar una sugerencia casual como una orden. Las decisiones importantes de rumbo deben pasar a una discusión formal, con opiniones, evidencia, responsabilidades y mecanismo de decisión explícitos. De lo contrario, el mayor peligro no es que el inversionista se equivoque, sino que el Founder malinterprete un comentario dicho de pasada.
\end{importantbox}

\subsection{Resumen de la sección}

La primera empresa le dejó a Xiao Hong tres lecciones de fondo: en un producto de herramienta, la comercialización debe empezar pronto; la financiación es una herramienta costosa, no la estrategia en sí, y el Founder debe conservar la soberanía de su juicio. Estas lecciones influyeron después directamente en sus decisiones sobre Monica, la expansión internacional y Manus.
